\subsection*{BR-010: further sound refinements of the permanent gate}
\textbf{Status: proved as a conservative finite-gate refinement.  Root whole-proof
and all six new-source/module reviews passed; no root rank rerun is claimed.
Only individually complete exclusions count; an unfinished search is
retained open.}

Two exact procedures refine the $596$ direct survivors of BR-009.
The combined remaining list is saved explicitly in
\texttt{certificates/per10\_combined\_gate.json}.  It removes the union
of the complete exclusions described below and the P-20 orbit from
BR-008.  The global minimum still lies in $\{17,18\}$ unless that list
becomes empty.  The final recorded gate has exactly $381$ remaining orbits, representing
$353304$ coordinate ten-subsets.  The completed conservative flag pass
excludes $210$ of the original $596$ cases, the all-endpoint substitution
pass excludes one further case, the arrangement postpass excludes three
further cases, and BR-008 excludes P-20.  These $215$ masks are distinct.
Both main passes processed all $596$ starting cases; the $386$ OPEN
flag outcomes remain inconclusive, including every timeout.  The
merger returned PASS.  These numerical counts are certificate results,
not assumptions in either proof.

\paragraph{Automatically certified connected solvable subgroups.}
For a fixed permanent coordinate ten-plane, let $Q_1,\ldots,Q_{10}$
be its literal four-entry slice matrices.  The ordered matrix
coordinates occurring in different $Q_a$ are disjoint.  For every
positive elementary row operation $r_j\mapsto r_j+s r_i$ with
$i<j$, and every positive elementary column operation, the module
constructor tries two actions: the same operation in $B,C$, and
opposite operations in $B,C$.  It first requires that the operation
preserve the used-variable spaces individually.  It then calculates
\[
 D Q_a=(D_B\otimes 1+1\otimes D_C)Q_a
\]
exactly, rejecting the operation unless every such matrix lies in
$\langle Q_1,\ldots,Q_{10}\rangle$.  Since the slice supports are
disjoint and every nonzero entry is one, this membership test is
particularly transparent: no coefficient outside the ten supports
is allowed, and the coefficients on each four-entry support must
all be equal.  Write the resulting matrix as
$DQ_a=\sum_b R_{ba}Q_b$.  The first-factor infinitesimal action is
$D_A(a_b)=-\sum_aR_{ba}a_a$.  The checker independently sums all three
factor actions on every tensor coefficient and asserts that the full
derivative is zero.

In characteristic zero the elementary unipotent actions integrate
these identities.  The induced action on the slice span is integrated
with its inverse dual action on $A$.  All $B,C$ matrices in the
construction are upper triangular, including the diagonal row/column
torus.  The subgroup preserving the slice span is a subgroup of a
product of upper-triangular groups, hence solvable.  The connected
algebraic subgroup generated by the torus and the retained
one-parameter unipotent groups is connected and solvable.  Its action
on $A$ is determined by its action on the concise slice span, so this
lift does not add a nonsolvable kernel.  The torus normalizes each
retained root.  No assumption about the full permanent stabilizer
or about the determinant Borel group is made.

If the tensor had border rank at most seventeen, the proper incidence
argument of BR-004 would supply a seventeen-plane $E\supset T(C^*)$
whose tautological tensor has border rank at most seventeen.  The
closed incidence is invariant under the constructed connected solvable
group, so a fixed such $E$ exists.  Lie's theorem supplies a complete
invariant flag in $E/T(C^*)$.  The exact joint-kernel enumerator
therefore lists every possibility.  A positive-dimensional kernel
is retained in its entirety.  The driver reuses the reviewed higher
Koszul, exact projective-line, projected-Koszul, and finite-continuation
engines.  A projective kernel of dimension above one is retained OPEN.
An unhandled nonrational continuation, an assertion failure, an error,
or a timeout likewise leaves the entire starting case OPEN.  Thus a
case receives EXCLUDED only if every rigid terminal extension and
every exported projective family has been excluded, including all
finite exceptional successors.  The per-case file retains the full
initial enumeration, all rigid tests, and the refinements indexed by
the original complete projective-family list.

\paragraph{All-endpoint border substitution.}
The second procedure is independent of that subgroup search.  Let
$D$ be either the ten-dimensional slice-label factor $A$ or the full
used-variable factor $B$, of dimension $d$.  Both are concise.  For
$1\le h<d$, multidimensional border substitution gives
\[
 \underline R(T)\ge h+
 \min_{L\in\operatorname{Gr}(d-h,D^*)}
 \underline R(T|_L).
\]
The stabilizing row/column torus has distinct weights in both $A$ and
$B$.  The rank-at-most-$r$ locus on this Grassmannian is closed and
torus invariant.  Therefore, if it is nonempty, it contains a torus
fixed point, which is a coordinate subspace.  Consequently the
minimum is achieved at one of all $\binom dh$ coordinate endpoints.
If the checker establishes
$\underline R(T|_L)\ge18-h$ for \emph{every} coordinate endpoint,
then $\underline R(T)\ge18$.

The script tries $h=1,2,3$ in $A$ and $h=1,2$ in $B$.  It reconstructs
each restriction literally and first tests its ordinary flattening.
For an $A$ restriction, all used-variable contraction vectors are
nonzero and have distinct weights, so their number is the flattening
rank.  For a $B$ restriction, zero $C$ contraction vectors are removed;
the remaining vectors still have distinct weights and are independent.
If necessary, exact modular Koszul ranks are then computed.  A map
with pure-tensor cost $c$ and rank greater than $(17-h)c$ proves the
required residual lower bound $18-h$.  The same integer-minor argument
as BR-009 lifts this rank witness from the prime field to
characteristic zero.  All endpoints of a successful attempt are saved
with their individual witnesses.  If even one endpoint has no sufficient
witness, the attempt is inconclusive.  A timeout is also inconclusive.
This procedure never adds $h$ to the lower bound of just one chosen
restriction.

\paragraph{The conservative union and its checker.}
The merger \texttt{merge\_per10\_refinements.py} validates every
claimed excluded case against the original complete $596$-case list.
For a flag exclusion, it checks that all terminal rigid flags appear
in the rigid-test list, that each failed a necessary condition, and
that the refinement indices are exactly all indices of the initial
projective-family list.  It checks that every refinement is excluded
and every finite continuation has no survivor or open family.  No
error-bearing record may count as an exclusion.

For a substitution exclusion it checks that the successful attempt
lists exactly all $\binom dh$ distinct coordinate endpoints in their
canonical order.  It recomputes the used dimensions and flattening
counts, recomputes each pure-tensor cost from its factor dimension,
and checks the strict numerical inequality at every endpoint.  The
ranks themselves were computed in the preceding source checker;
the merger explicitly records that it does not rerun those ranks.
It then removes the union of these validated masks and the proved
P-20 mask $2015$.  A processed-but-open or wholly unprocessed case
remains in the output list.  Hence partial phase completion cannot
make the combined gate unsound.

The scripts are \texttt{permanent\_solvable.py},
\texttt{refine\_per10\_global\_gate.py},
\texttt{check\_per10\_substitution.py}, and
\texttt{merge\_per10\_refinements.py}.  The first and third procedures
and the merger use standard-library exact arithmetic; the projective
pencil refinement additionally uses SymPy.  Runtime limits are only
search limits and never algebraic exclusions.  A restart resumes
only completed per-case records, preserving all OPEN statuses.

\paragraph{A small exact arrangement refinement.}
The optional postpass \texttt{check\_per10\_arrangements.py}, with
\texttt{first\_wedge\_arrangement.py}, also treats higher-dimensional
one-weight extension families conservatively.  Let $E_0$ be a partial
invariant flag and let $v$ vary in one weight-space kernel $K$.
For the first Koszul map along $A$, the additional columns are
$v\wedge a_i$.  Their output weights are the weight of $v$ plus the
distinct weights of the $a_i$, hence are pairwise distinct.  After
quotienting by the image of the map on $E_0$, each new column is either
zero or adds exactly one to the rank, independently of the others.
The same statement holds along $B$.  Each zero condition is the kernel
of an explicitly computed linear map on $K$.

The checker enumerates every intersection of these linear kernels by
canonical rational row reduction, retaining every nonzero parameter
subspace.  On the generic stratum of such an intersection, the set of
zero columns is exactly the set of kernel conditions containing the
whole intersection.  This stratum is nonempty whenever needed, since
the field is infinite and finitely many proper linear subspaces cannot
cover a vector space.  More directly, for any candidate vector, the
intersection of all its zero conditions appears in the enumeration;
its generic nonzero-column count is the same as that of the candidate.
Thus every candidate is covered.  Adding the number of nonzero
columns to the two base ranks gives the exact first-Koszul ranks on
each zero pattern.  Intersections whose generic ranks satisfy both
necessary thresholds may be retained in their entirety, because
further specialization only decreases those ranks.

The postpass uses this result only when no allowed nonzero parameter
subspace remains.  Any nonempty remaining parameter space stays OPEN.
A complete initial flag cover already saved by the preceding procedure
is required; all its terminal rigid tests must have failed, and each
projective family must either have its previous complete exclusion or
receive this empty-arrangement certificate.  A timeout stays OPEN.
The merger checks this complete index cover before adding the resulting
masks to the exclusion union.  This postpass is the same linear
zero-pattern method used in BR-001, now applied to a general invariant
partial flag; no sampled parameter or generic-rank-only inference is
used.
